\section{Why Responsible AI}

\begin{frame}{Where This Lecture Sits}

    So far the course has optimised a single number: accuracy, F1, RMSE. That number is computed
    \textbf{over the whole test set}.

    \vspace{1em}

    \begin{block}{The gap this lecture closes}
        A model that is accurate \emph{on average} can still be systematically wrong for a
        particular group of people, and can be right for reasons you would not accept if you
        could see them.
    \end{block}

    \vspace{1em}

    \begin{itemize}
        \setlength\itemsep{0.5em}
        \item \textbf{Part 1 --- Fairness:} how to find and reduce systematic error across groups.
        \item \textbf{Part 2 --- Interpretability:} how to find out what a model is actually using.
        \item \textbf{Part 3 --- Documentation and governance:} how to write it down and who checks it.
    \end{itemize}

    \vspace{0.5em}
    \centering
    \small
    \textit{Not covered here:} agent- and LLM-specific safety --- prompt injection, tool misuse,
    excessive agency --- which is the subject of a dedicated lecture in \textbf{the Advanced NLP track (Phase 2B)}.

    \vspace{0.6em}
    \normalsize
    \textit{We start with four documented cases. Each one failed in a different place in the pipeline.}
\end{frame}
